\section{VulnRAG: System Architecture}
\label{sec:design}

VulnRAG is designed as a stateful, modular framework that decouples vulnerability reasoning from raw data retrieval. Unlike traditional monolithic RAG pipelines, our architecture (Fig.~\ref{fig:arch}) utilizes a Master-Worker hierarchy to coordinate specialized agents across heterogeneous LLM backends.

\subsection{Knowledge Synchronization and Storage}
The foundation of VulnRAG is a continuously updated knowledge base of vulnerability disclosures. 
\begin{itemize}
    \item \textbf{Continuous Sync}: The `nvd_sync_service.py` module manages real-time synchronization with the NIST NVD API 2.0. It parses raw JSON advisories into a structured format, normalizing software versioning (CPE) and security metrics (CVSS).
    \item \textbf{Vector Store}: We utilize **Qdrant** for high-performance vector storage. CVE descriptions and metadata are embedded using the \textbf{BAAI/bge-large-en-v1.5} model via the `SentenceTransformers` framework.
    \item \textbf{Metadata Filtering}: To improve retrieval precision, we implement hybrid search patterns in `cve_vector_store.py` that combine semantic embeddings with payload-based metadata filters (e.g., CVSS severity, vendor tags, and exploit availability).
\end{itemize}

\subsection{Master-Worker Orchestration}
The system's reasoning logic is managed by the `VulnIntelAgent`, which orchestrates the following specialized sub-modules:
\begin{enumerate}
    \item \textbf{Query Reformulator}: Enhances raw user queries with domain-specific terminology.
    \item \textbf{Smart Searcher}: Coordinates the retrieval and reranking loop.
    \item \textbf{Graph Explorer}: Discovers structural relationships between vulnerabilities.
    \item \textbf{Evaluator}: Assesses result quality and triggers self-correction loops.
\end{enumerate}

By modularizing these components, VulnRAG can adapt its reasoning strategy based on the complexity of the security task at hand.
